\section{Multivariate Continuous Random Variables}

\subsection{Distribution, Density and Independence}

\begin{definition}[Distribution and Density of Joint Continuous Random Variable]
    \label{def:distribution-density-joint-crv}%
    The joint continuous random variable \(\vb{X} = \pqty{X_1, \ldots, X_n} \in \RR^n\) has \emph{probability density function} \(f\), if for all \(x_1, \ldots, x_n\), that
    \[
        \Prob\pqty{X_1 \leq x_1, \ldots, X_n \leq x_n} = \int_{\minusinfty}^{x_1} \dd{y_1} \cdots \int_{\minusinfty}^{x_n} \dd{y_n} f\pqty{y_1, \ldots, y_n}.
    \]

    The left-hand side
    \[
        F \pqty{x_1, \ldots, x_n} = \Prob\pqty{X_1 \leq x_1, \ldots, X_n \leq x_n}
    \]
    is the \emph{probability distribution function} of \(\vb{X}\).
\end{definition}

We have
\[
    f\pqty{x_1, \ldots, x_n} = \frac{\partial^n F}{\partial x_1 \cdots \partial x_n} \pqty{x_1, \ldots, x_n}.
\]

For all \(\mathcal{B} \subseteq \RR^n\), we have
\[
    \Prob\pqty{\pqty{X_1, \ldots, X_n} \in \mathcal{B}} = \int_{\mathcal{B}} f \pqty{y_1, \ldots, y_n} \dd[n]{\vb{y}}.
\]

\begin{definition}[Independence of Continuous Random Vaiables]
    \label{def:independence-continuous-random-variables}%
    We say the continuous random variables \(X_1, \ldots, X_n\) are \emph{independent}, if for all \(x_1, \ldots, x_n\), we have
    \[
        \Prob\pqty{X_1 \leq x_1, \ldots, X_n \leq x_n} = \Prob\pqty{X_1 \leq x_1} \cdots \Prob\pqty{X_n \leq x_n}.
    \]
\end{definition}

\begin{theorem}[Independence if and only if Density is Separable]
    \label{thm:rv-independent-density-separable}%
    Let \(\vb{X} = \pqty{X_1, \ldots, X_n}\) have density \(f\).
    \begin{enumerate}
        \item Suppose \(X_1, \ldots, X_n\) are independent, with densities \(f_1, \ldots, f_n\) respectively. Then
        \[
            f\pqty{x_1, \ldots, x_n} = f_1\pqty{x_1} \cdots f_n\pqty{x_n}
        \]
        for all \(x_1, \ldots, x_n\).
        \item Suppose \(f\) factorises into above, for some non-negative functions \(f_1, \ldots, f_n\). Then \(X_1, \ldots, X_n\) are independent, with densities proportional to \(f_1, \ldots, f_n\) respectively (i.e. up to normalisation).
    \end{enumerate}
\end{theorem}

\begin{proof}[of First Part]
    We investigate the distribution function. We have
    \begin{align*}
        \Prob\pqty{X_1 \leq x_1, \ldots, X_n \leq x_n} &= \Prob\pqty{X_1 \leq x_1} \cdots \Prob\pqty{X_n \leq x_n} \\
        &= \int_{\minusinfty}^{x_1} f_1 \pqty{y_1} \dd{y_1} \cdots \int_{\minusinfty}^{x_n} f_n \pqty{y_n} \dd{y_n} \\
        &= \int_{\minusinfty}^{x_1} \dd{y_1} \cdots \int_{\minusinfty}^{x_n} \dd{y_n} f_1\pqty{y_1} f_2\pqty{y_2} \cdots f_n\pqty{y_n},
    \end{align*}
    so
    \[
        f\pqty{x_1, \ldots, x_n} = f_1\pqty{x_1} \cdots f_n\pqty{x_n}.
    \]
\end{proof}

\begin{proof}[of Second Part]
    We have that
    \begin{align*}
        \Prob\pqty{X_1 \leq x_1, \ldots, X_n \leq x_n} &= \int_{\minusinfty}^{x_1} \dd{y_1} \cdots \int_{\minusinfty}^{x_n} \dd{y_n} f\pqty{y_1, \ldots, y_n} \\
        &= \int_{\minusinfty}^{x_1} \dd{y_1} \cdots \int_{\minusinfty}^{x_n} \dd{y_n} f_1\pqty{y_1} \cdots f_n\pqty{y_n} \\
        &= \int_{\minusinfty}^{x_1} f_1\pqty{y_1} \dd{y_1} \cdots \int_{\minusinfty}^{x_n} f_n\pqty{y_n} \dd{y_n}.
    \end{align*}

    But we have that
    \[
        \int_{\minusinfty}^{\plusinfty} \dd{x_1} \cdots \int_{\minusinfty}^{\plusinfty} \dd{x_n} f_1\pqty{x_1} \cdots f_n\pqty{x_2} = 1,
    \]
    and hence
    \[
        \Prob\pqty{X_1 \leq x_1, \ldots, X_n \leq x_n} = \frac{\int_{\minusinfty}^{x_1} f_1\pqty{y_1} \dd{y_1}}{\int_{\minusinfty}^{\plusinfty} f_1\pqty{y_1} \dd{y_1}} \cdots \frac{\int_{\minusinfty}^{x_1} f_n\pqty{y_n} \dd{y_n}}{\int_{\minusinfty}^{\plusinfty} f_n\pqty{y_n} \dd{y_n}}.
    \]

    So \(X_1, \ldots, X_n\) are independent, with densities
    \[
        \frac{f_i}{\int_{\minusinfty}^{\plusinfty} f_i\pqty{x} \dd{x}}
    \]
    respectively.
\end{proof}

\subsection{Marginal Distribution and Density}

\begin{definition}[Marginal Distribution and Density]
    \label{def:marginal-distribution-density}%
    For a joint continuous random variable \(\vb{X} = \pqty{X_1, X_2, \ldots, X_n}\), the \emph{marginal distribution} \(F_{X_i}\) and the \emph{marginal density} \(f_{X_i}\) of \(X_i\) are given by, respectively,
    \[
        F_{X_i} \pqty{x} = \Prob\pqty{X_i \leq x} = \Prob\pqty{X_1 \in \RR, \ldots, X_{i - 1} \in \RR, X_i \leq x, X_{i + 1} \in \RR, \ldots, X_n \in \RR}
    \]
    and
    \[
        f_{X_i} \pqty{x} = F_{X_i}'\pqty{x}.
    \]
\end{definition}

\begin{proposition}[Integration Formula for Marginal Density and Distribution]
    \label{prop:marginal-distribution-density-formula}%
    Suppose \(\vb{X} = \pqty{X_1, \ldots, X_n}\) has density \(f\). Then we have
    \[
        F_{X_1} \pqty{x} = \int_{\minusinfty}^{x} \dd{x_1} \cdot \int_{\minusinfty}^{\plusinfty} \dd{x_2} \cdots \int_{\minusinfty}^{\plusinfty} \dd{x_n} f\pqty{x_1, \ldots, x_n},
    \]
    and
    \[
        f_{X_1} \pqty{x} = \int_{\minusinfty}^{\plusinfty} \dd{x_2} \cdots \int_{\minusinfty}^{\plusinfty} \dd{x_n} f\pqty{x, \ldots, x_n}.
    \]
\end{proposition}

\subsection{Sum of Independent Random Variables}

Suppose \(X\) and \(Y\) are independent random variables with densities \(f_X\) and \(f_Y\) respectively. What is the density of \(X + Y\)?

In the discrete case, we recall this is given by the convolution
\[
    \Prob\pqty{X + Y = z} = \sum_{x = \minusinfty}^{\plusinfty} \Prob\pqty{X = x} \cdot \Prob\pqty{Y = z - x}.
\]

The continuous case is similar.

\begin{proposition}[Distribution of Sum of Independent Continuous Random Variables]
    \label{prop:dist-sum-indp-crv}%
    Let \(X, Y\) be independent, continuous random variables. For the sum, we have
    \[
        f_{X + Y}\pqty{z} = \int_{\minusinfty}^{\infty} f_X\pqty{x} f_Y\pqty{z - x} \dd{x}.
    \]
\end{proposition}

\begin{proof}
    Investigating the distribution function, we have
    \begin{align*}
        \Prob\pqty{X + Y \leq z} &= \Prob\pqty{\pqty{X, Y} \in \set{\pqty{x, y} \in \RR^2}{x + y \leq z}} \\
        &= \iint_{x + y \leq z} \dd[2]{\pqty{x, y}} f_{X, Y} \pqty{x, y} \\
        &= \iint_{x + y \leq z} \dd[2]{\pqty{x, y}} f_{X} \pqty{x} f_{Y} \pqty{y} \\
        &= \int_{\minusinfty}^{\plusinfty} \dd{x} \int_{\minusinfty}^{z - x} \dd{y} f_X\pqty{x} f_Y\pqty{y} \\
        &= \int_{\minusinfty}^{\plusinfty} \dd{x} \int_{\minusinfty}^{z} \dd{y} f_X\pqty{x} f_Y\pqty{y - z} \\
        &= \int_{\minusinfty}^{z} \dd{y} \pqty{\int_{\minusinfty}^{\plusinfty} f_X\pqty{x} f_Y\pqty{y - x} \dd{x}}.
    \end{align*}

    Differentiating this gives the result, as desired.
\end{proof}

\begin{definition}[Convolution of Continuous Random Variables]
    \label{def:convolution-crv}%
    Let \(X, Y\) be two continuous random variables. Then the distribution of \(X + Y\), given by
    \[
        \int_{\minusinfty}^{\infty} f_X\pqty{x} f_Y\pqty{z - x} \dd{x}
    \]
    is known as the \emph{convolution} of the two random variables.
\end{definition}

\subsection{Conditional Density}

\begin{definition}[Conditional Density]
    \label{def:conditional-density}%
    Let \(X, Y\) be two continuous random variables, with joint density \(f_{X, Y}\) and marginal densities \(f_{X}\) and \(f_{Y}\). The \emph{conditional density} of \(X\) given \(Y = y\) is
    \[
        f_{\bound{\Conditional{X}{Y}}} \pqty{\Conditional{x}{y}} = \frac{f_{X, Y}\pqty{x, y}}{f_Y\pqty{y}}
    \]
    for \(f_Y\pqty{y} > 0\).
\end{definition}

\begin{theorem}[Continuous Version of Law of Total Probability]
    \label{thm:cont-law-total-prob}%
    For continuous random variables \(X\) and \(Y\), we have
    \[
        f_X\pqty{x} = \int_{\minusinfty}^{\plusinfty} f_{X, Y} \pqty{x, y} \dd{y} = \int_{\minusinfty}^{\plusinfty} f_{\bound{\Conditional{X}{Y}}} \pqty{\Conditional{x}{y}} f_Y\pqty{y} \dd{y}.
    \]
\end{theorem}

\begin{definition}[Conditional Expectation for Continuous Random Variable]
    \label{def:conditional-expectation-crv}%
    For continuous random variables \(X\) and \(Y\), the \emph{conditional expectation} of \(X\) given \(Y\) is defined as
    \[
        \Expt\bqty{\Conditional{X}{Y}} = g\pqty{Y},
    \]
    where we have
    \[
        g\pqty{y} = \Expt\bqty{\Conditional{X}{Y = y}} = \int_{\minusinfty}^{\plusinfty} x f_{\bound{\Conditional{X}{Y}}} \pqty{\Conditional{x}{y}} \dd{x}.
    \]
\end{definition}

\subsection{Transformation of Continuous Random Variables}

\begin{theorem}[Transformation of Joint Continuous Random Variables]
    \label{thm:transformation-joint-continuous-random-variable}%
    Let \(\vb{X}\) be a random variable with values in \(\mathcal{D} \subseteq \RR^n\) and density \(f\). Let
    \[
        \functiondomain{g}{\mathcal{D}}{g\pqty{\mathcal{D}} \subseteq \RR^n}
    \]
    be a bijection, with a continuous derivative on \(\mathcal{D}\), with
    \[
        \det g'\pqty{\vb{x}} = \det \pqty{\pdv{g\pqty{\vb{x}}_i}{x_j}} \neq 0
    \]
    for any \(\vb{x} \in \mathcal{D}\). Then for any \(\vb{y} \in g\pqty{\mathcal{D}}\), the random variable \(\vb{Y} = g\pqty{\vb{X}}\) has density
    \[
        f_{\vb{Y}}\pqty{\vb{y}} = f_{\vb{X}} \pqty{g\inverse\pqty{\vb{y}}} \cdot \abs{\det \vb{J}},
    \]
    where \(\vb{J}\) is the \emph{Jacobian}, given by
    \[
        J_{i, j} = \pdv{x_i}{y_j}
    \]
    where \(\vb{y} = g\pqty{\vb{x}}\).
\end{theorem}

\begin{example}[Transformation of Cartesian and Polar Form]
    \label{ex:transformation-cartesian-polar}%
    Let \(X, Y \sim \Normal\pqty{0, 1}\) be independent, and represent Cartesian coordinates in the plane. Let \(R, \Theta\) be the random variables for the usual polar coordinates, i.e.,
    \[
        X = R \cos \Theta, \quad Y = R \sin \Theta,
    \]
    restricting us to \(\Theta \in \cointv{0}{2\pi}\).

    Applying \autoref{thm:transformation-joint-continuous-random-variable}, we have
    \[
        f_{\pqty{R, \Theta}} \pqty{r, \theta} = f_{\pqty{X, Y}} \pqty{r \cos \theta, r \sin \theta} \cdot \abs{\det \vb{J}},
    \]
    where the Jacobian is given by
    \[
        \vb{J} = \pmqty{
            \cos \theta & - r \sin \theta \\
            \sin \theta & r \cos \theta
        },
    \]
    and so \(\abs{\det \vb{J}} = r\). Hence, using independence,
    \begin{align*}
        f_{\pqty{R, \Theta}} \pqty{r, \theta} &= f_X \pqty{r \cos \theta} \cdot f_Y \pqty{r \sin \theta} \cdot r \\
        &= r \cdot \frac{1}{\sqrt{2\pi}} \cdot e^{- \frac{r^2 \cos^2 \theta}{2}} \cdot \frac{1}{\sqrt{2 \pi}} \cdot e^{- \frac{r^2 \sin^2 \theta}{2}} \\
        &= \frac{1}{2\pi} \cdot r \cdot e^{- \frac{r^2}{2}}.
    \end{align*}

    Therefore, \(\Theta \sim \Uniform\ccintv{0}{2\pi}\), and \(R \sim r e^{-\frac{r^2}{2}}\) for \(r > 0\). They are independent.
\end{example}

\subsection{Order Statistics of a Random Variable}

Let \(X_1, \ldots, X_n\) be independent and identically distributed random variables, with probability distribution function \(F\) and density function \(f\).

We order them from smallest to largest, which then gives
\[
    X_{\pqty{1}} \leq X_{\pqty{2}} \leq \cdots \leq X_{\pqty{n}},
\]
and we let \(Y_i = X_{\pqty{i}}\).

\begin{definition}[Order Statistics]
    \label{def:order-statistics}%
    The \(Y_i\)s found above are the \emph{order statistics} of the sample.
\end{definition}

\begin{example}[Density of \(Y_1\) and \(Y_n\)]
    \label{ex:density-y1-yn}%
    For \(Y_1\), we have
    \[
        \Prob\pqty{Y_1 \leq x} = 1 - \Prob\pqty{Y_1 > x} = 1 - \pqty{\Prob\pqty{X_1 > x}}^n = 1 - \pqty{1 - F\pqty{x}}^n,
    \]
    and hence
    \[
        f_{Y_1} \pqty{x} = \dv{x} \Prob\pqty{Y_1 \leq x} = n \pqty{1 - F\pqty{x}}^{n - 1} f\pqty{x}.
    \]

    From a similar argument, we have
    \[
        \Prob\pqty{Y_n \leq x} = \Prob\pqty{X_1 \leq x}^n = \pqty{F\pqty{x}}^n.
    \]

    Hence, differentiating gives
    \[
        f_{Y_n} \pqty{x} = n \pqty{F\pqty{x}}^{n - 1} f\pqty{x}.
\]
\end{example}

\begin{proposition}[Density of Order Statistics]
    \label{prop:density-order-stats}%
    We have
    \[
        f_{\pqty{Y_1, \ldots, Y_n}} \pqty{x_1, \ldots, x_n} = \begin{cases}
            n! f\pqty{x_1} \cdots f\pqty{x_n}, & x_1 < \cdots < x_n, \\
            0, & \text{otherwise}.
        \end{cases}
    \]
\end{proposition}

\begin{proof}
    For \(x_1 < x_2 < \cdots < x_n\), we have
    \begin{align*}
        F_{\pqty{Y_1, \ldots, Y_n}} \pqty{x_1, \ldots, x_n} &= \Prob\pqty{Y_1 \leq x_1, \ldots, Y_n \leq x_n} \\
        &= n! \cdot \Prob\pqty{X_1 \leq x_1, \ldots, X_n \leq x_n, X_1 \leq \cdots \leq X_n} \\
        &= n! \cdot \int_{\RR} \dd{u_1} \cdots \int_{\RR} \dd{u_n} \cdot f_{\pqty{X_1, \ldots, X_n}} \pqty{u_1, \ldots, u_n} \cdot \\
        &\phantom{=} \quad \indicator \pqty{u_1 \leq x_1, \ldots, u_n \leq x_n, u_1 \leq \cdots \leq u_n} \\
        &= n! \cdot \int_{\minusinfty}^{x_1} \dd{u_1} \int_{u_1}^{x_2} \dd{u_2} \cdots \int_{u_{n - 1}}^{x_n} \dd{u_n} \cdot f\pqty{u_1} \cdots f\pqty{u_n}.
    \end{align*}

    Differentiating gives us
    \[
        f_{\pqty{Y_1, \ldots, Y_n}} \pqty{x_1, \ldots, x_n} = \begin{cases}
            n! \cdot f\pqty{x_1} \cdots f\pqty{x_n}, & x_1 < \cdots < x_n, \\
            0, & \text{otherwise}.
        \end{cases}
    \]
\end{proof}

In particular, note from the above that the \(Y_i\)s are not independent -- although this density function seems to be separable, it is piecewise where the pieces depend on the relation between the variables.

\begin{example}[Order Statistics for Exponentials]
    \label{ex:order-statistics-exp}%
    Let \(X \sim \Exponential \pqty{\lambda}\) and \(Y \sim \Exponential \pqty{\mu}\), with \(X \indp Y\) and \(\lambda, \mu > 0\). Set \(Z = \min\Bqty{X, Y}\).

    We have for \(z > 0\), that
    \[
        \Prob\pqty{Z \leq z} = 1 - \Prob\pqty{Z \geq z} = 1 - \Prob\pqty{X \geq z} \cdot \Prob\pqty{Y \geq z} = 1 - e^{- \pqty{\lambda + \mu} z}.
    \]

    This gives \(Z \sim \Exponential\pqty{\lambda + \mu}\).

    Similarly, if \(X_i \sim \Exponential \pqty{\lambda_i}\) are independent, we have
    \[
        \min\Bqty{X_1, \ldots, X_n} \sim \Exponential\pqty{\sum_{i = 1}^{n} \lambda_i}.
    \]

    Now, suppose \(\lambda_1 = \cdots = \lambda_n = \lambda\), and let \(Y_1, Y_2, \ldots, Y_n\) be the order statistics.

    We \(Z_1 = Y_1, Z_2 = Y_2 - Y_1, \ldots, Z_i = Y_i - Y_{i - 1}\). We find the density of \(\pqty{Z_1, \ldots, Z_n}\).

    Let
    \[
        \vb{Z} = \pmqty{Z_1 \\ \vdots \\ Z_n},
        \vb{Y} = \pmqty{Y_1 \\ \vdots \\ Y_n},
        \vb{A} = \pmqty{
             1 &  0 & 0 & \cdots &  0 &  0 &  0 \\
            -1 &  1 & 0 & \cdots &  0 &  0 &  0 \\
             0 & -1 & 1 & \cdots &  0 &  0 &  0 \\
             \vdots & \vdots & \vdots & \ddots & \vdots & \vdots & \vdots \\
             0 &  0 & 0 & \cdots & -1 &  1 & 0 \\
             0 &  0 & 0 & \cdots &  0 & -1 & 1
        },
    \]
    so we have \(\vb{Z} = \vb{A} \vb{y}\), and \(\det \vb{A} = 1\). Hence, let \(\vb{z} = \vb{A} \vb{y}\), we have
    \begin{align*}
        f_{\pqty{Z_1, \ldots, Z_n}} \pqty{z_1, \ldots, z_n} &= f_{\pqty{Y_1, \ldots, Y_n}} \pqty{y_1, \ldots, y_n} \abs{\det \vb{A}} \\
        &= n! \cdot f\pqty{y_1} \cdots f\pqty{y_n} \\
        &= n! \lambda e^{-\lambda y_1} \cdots \lambda e^{-\lambda y_n} \\
        &= \prod_{i = 1}^{n} \pqty{\pqty{n - i + 1} \lambda e^{-\lambda \pqty{n - i + 1} z_i}}.
    \end{align*}

    So the \(Z_i\)s are independent, and
    \[
        Z_i \sim \Exponential\pqty{\lambda \pqty{n - i + 1}}.
    \]
\end{example}

\subsection{Expectation and Variance Algebra}

\begin{definition}[Expectation and Variance of Multivariate Random Variable]
    \label{def:expt-var-multivariate}%
    For a random variable \(\vb{X}\), we define the \emph{mean}
    \[
        \vb*{\mu} = \Expt\bqty{\vb{X}} = \pmqty{\Expt\bqty{X_1} \\ \vdots \\ \Expt\bqty{X_n}},
    \]
    and the \emph{variance}
    \[
        \Var\pqty{\vb{X}} = \Expt\bqty{\pqty{\vb{X} - \vb*{\mu}} \pqty{\vb{X} - \vb*{\mu}}\transpose}.
    \]

    Equivalently, the variance is defined as
    \[
        \Var\pqty{\vb{X}} = \pqty{\Cov\pqty{X_i, X_j}}_{i, j}.
    \]
\end{definition}

\begin{proposition}[Expectation and Variance of Multivariate Random Variable]
    \label{prop:expt-var-transform-multivariate-rv}%
    We have
    \[
        \Expt\bqty{\vb{u} \transpose \vb{X}} = \vb{u}\transpose \Expt\bqty{\vb{X}}
    \]
    and
    \[
        \Var\pqty{\vb{u} \transpose \vb{X}} = \vb{u}\transpose \Var\pqty{\vb{X}} \vb{u}.
    \]
\end{proposition}

\begin{proof}
    \[
        \Expt\bqty{\vb{u} \transpose \vb{X}} = \Expt\bqty{\sum_{i = 1}^{n} u_i X_i} = \sum_{i = 1}^{n} u_i \Expt\bqty{X_i} = \vb{u}\transpose \Expt\bqty{X},
    \]
    and
    \[
        \Var\pqty{\vb{u} \transpose \vb{X}} = \Var\pqty{\sum_{i = 1}^{n} u_i X_i} = \sum_{i, j = 1}^{n} u_i \Cov\pqty{X_i, X_j} u_j = \vb{u}\transpose \vb{\Var\pqty{\vb{X}}} \vb{u}. 
    \]
\end{proof}

\begin{definition}[Positive/Negative Definiteness]
    \label{def:positive-negative-definiteness}%
    Let \(\vb{M}\) be an \(n \times n\) real symmetric matrix. Then we say it is
    \begin{itemize}
        \item \emph{positive definite}, denoted as \(\vb{M} \posDef 0\), if for all \(\vb{u} \in \RR^n \setminus \Bqty{\vb{0}}\), we have \(\vb{u}\transpose \vb{M} \vb{u} > 0\);
        \item \emph{positive semi-definite}, or \emph{non-negative definite}, denoted as \(\vb{M} \posDefEq 0\), if for all \(\vb{u} \in \RR^n\), we have \(\vb{u} \transpose \vb{M} \vb{u} \geq 0\);
        \item \emph{negative definite}, denoted as \(\vb{M} \negDef 0\), if for all \(\vb{u} \in \RR^n \setminus \Bqty{\vb{0}}\), we have \(\vb{u}\transpose \vb{M} \vb{u} < 0\);
        \item \emph{negative semi-definite}, or \emph{non-positive definite}, denoted as \(\vb{M} \negDefEq 0\), if for all \(\vb{u} \in \RR^n\), we have \(\vb{u} \transpose \vb{M} \vb{u} \leq 0\).
    \end{itemize}

    Equivalently, in terms of eigenvalues, it is
    \begin{itemize}
        \item \emph{positive definite}, if all its eigenvalues are positive;
        \item \emph{positive semi-definite}, or \emph{non-negative definite}, if all its eigenvalues are non-negative;
        \item \emph{negative definite}, if all its eigenvalues are negative;
        \item \emph{negative semi-definite}, or \emph{non-positive definite}, if all its eigenvalues are non-positive.
    \end{itemize}
\end{definition}

\begin{claim}[Variance is Non-Negative Definite]
    \label{claim:variance-non-negative-definite}%
    The variance of a multivariate distribution is non-negative definite.
\end{claim}

\begin{proof}
    We have
    \[
        \Var\pqty{\vb{u}\transpose \vb{X}} = \vb{u}\transpose \Var\pqty{\vb{X}} \vb{u} \geq 0.
    \]
\end{proof}

\begin{definition}[Square Root of Non-Negative Definite Matrix]
    \label{def:sqrt-non-negative-definite-matrix}%
    Let \(\vb{V}\) be a non-negative definite matrix, and let its orthogonal decomposition be
    \[
        \vb{V} = \vb{U}\transpose \vb{D} \vb{U}
    \]
    for an orthogonal matrix \(\vb{U}\transpose = \vb{U}\inverse\), and \(\vb{D}\) diagonal with
    \[
        \vb{D} = \diag\pqty{\lambda_1, \ldots, \lambda_n} = \pmqty{\dmat{\lambda_1, \ddots, \lambda_n}}
    \]
    with \(\lambda_i \geq 0\) for \(i = 1, \ldots, n\).

    The \emph{square root} of \(\vb{V}\) is defined as
    \[
        \vb*{\sigma} = \vb{U}\transpose \sqrt{\vb{D}} \vb{U}
    \]
    where
    \[
        \sqrt{\vb{D}} = \diag\pqty{\sqrt{\lambda_1}, \ldots, \sqrt{\lambda_n}} = \pmqty{\dmat{\sqrt{\lambda_1}, \ddots, \sqrt{\lambda_n}}}
    \]
\end{definition}

\begin{claim}[Independence implies Variance is Diagonal]
    \label{claim:indep-diagonal-variance}%
    If \(X_1, \ldots, X_n\) are independent, then \(\vb{V}\) is a diagonal matrix.
\end{claim}

\begin{proof}
    The non-diagonal arguments are pairwise covariances between \(X_i\) and \(X_j\) where \(i \neq j\).
\end{proof}

\subsection{Correlation}

We restrict us to the two-dimensional case. Set \(\mu_1 = \Expt\bqty{X_1}, \mu_2 = \Expt\bqty{X_2}, \sigma_1^2 = \Var\pqty{X_1}, \sigma_2^2 = \Var\pqty{X_2}\).

\begin{definition}[Correlation]
   \label{def:corr}%
    For two random variables \(X_1\) and \(X_2\), their \emph{correlation} is defined as
    \[
        \rho = \Corr\pqty{X_1, X_2} = \frac{\Cov\pqty{X_1, X_2}}{\sqrt{\Var\pqty{X_1} \Var\pqty{X_2}}}.
    \]
\end{definition}

\begin{claim}[Correlation has Absolute Value no more than \(1\)]
    \label{claim:corr-abs-val-no-more-one}%
    The correlation satisfies \(\rho \in \ccintv{-1}{1}\).
\end{claim}

\begin{proof}
    This simply comes from Cauchy--Schwartz inequality.
\end{proof}

\begin{claim}[Variance Matrix in terms of Correlation]
    \label{claim:var-mat-corr}%
    Let \(\sigma_1, \sigma_2 > 0\) and \(\rho \in \ccintv{-1}{1}\). Then the matrix
    \[
        \vb{V} = \pmqty{
            \sigma_1^2 & \rho \sigma_1 \sigma_2 \\
            \rho \sigma_1 \sigma_2 & \sigma_2^2
        }
    \]
    is non-negative definite.
\end{claim}

\begin{proof}
    Let \(\vb{u} \in \RR^2\) be arbitrary, and let \(\vb{u} = \pmqty{u_1 \\ u_2}\).

    We have
    \begin{align*}
        \vb{u}\transpose \vb{V} \vb{u} &= \pqty{1 - \rho} \pqty{\sigma_1^2 u_1^2 + \sigma_2^2 u_2^2} + \rho \pqty{\sigma_1 u_1 + \sigma_2 u_2}^2 \\
        &= \pqty{1 + \rho} \pqty{\sigma_1^2 u_1^2 + \sigma_2^2 u_2^2} - \rho \pqty{\sigma_1 u_1 - \sigma_2 u_2}^2.
    \end{align*}

    If \(\rho \in \ccintv{-1}{0}\) we use the second equality, and if \(\rho \in \ccintv{0}{1}\) we use the first equality.
\end{proof}